\chapter{Phonology: Empirically Derived Invariants}

\section{Transition from Theoretical to Empirical Modeling}
Previous iterations of the PED framework utilized idealized sound laws which failed to account for the full variance of the primary attested data. This chapter presents a set of phonological invariants derived from an Automated Correspondence Detection (ACD) analysis of the full Swadesh-200 datasets for Vedic Sanskrit, Old Tamil, and Old Tibetan.

\section{The Radial Nasal Shift (*N)}
The most widespread signal in the Eurasian substrate is the tripartite divergence of the ancestral nasal.
\begin{itemize}
    \item \textbf{Western Node (Vedic):} Shift to bilabial \textit{*m} (e.g., \textit{*nga > aham}, \textit{*nam > vayam}).
    \item \textbf{Southern Node (Tamil):} Retention of alveolar \textit{*n} or palatal \textit{*ñ} (e.g., \textit{*nga > yāṉ}, \textit{*nyi > ñāyiṟu}).
    \item \textbf{Eastern Node (Tibetan):} Retention of velar \textit{*ng} or palatal \textit{*ny} (e.g., \textit{*nga > nga}, \textit{*nyi > nyi-ma}).
\end{itemize}

\section{The Dental-Palatal Invariant (*T)}
The ancestral dental invariant governs core topography and relational mathematics. 
\begin{center}
\textit{*T $\rightarrow$ d (Vedic), t/m (Tamil), sh/ny (Tibetan)}
\end{center}
The Tibetan node exhibits characteristic palatalization (\textit{*T > sh/ny}) before high vowels, while the Southern node often exhibits nasalization of the onset in deictic and numeral categories.

\section{Labial and Velar Radiation Patterns}
Beyond the nasal and dental cores, we identify two primary radiation patterns governing the ancestral stop series:
\begin{itemize}
    \item \textbf{Labial Radiation (*P):} Diverges into Western voiced \textit{b}, Southern voiceless \textit{p}, and Eastern aspirated \textit{ph} (e.g., \textit{*pa-r > paru/peru/lci}).
    \item \textbf{Velar Radiation (*K):} Diverges into Western \textit{k/g}, Southern \textit{k/c} (via palatalization), and Eastern \textit{kh/g} (e.g., \textit{*ka-n > karṇa/cevi/rna}).
\end{itemize}

\section{Sibilant Palatalization and Liquid Drift}
\begin{itemize}
    \item \textbf{Sibilant Shift (*S):} The ancestral sibilant frequently palatalizes in the Eastern and Southern nodes (e.g., \textit{*sa-n > māṃsa/ūṇ/sha}).
    \item \textbf{Liquid Drift (*R/L):} Liquids exhibit high stability but frequently interchange between alveolar and retroflex positions in the Southern node.
\end{itemize}

\section{Conclusion: The Empirical Sound Law Engine}
By integrating these tripartite radiation patterns, we establish a systematic framework for deriving the attested Eurasian lexicon. This empirical engine removes the need for 'idealized' intermediate forms, grounding the PED reconstruction in documented phonetic history.

